\section{References}

\subsection*{Sacred Scripture}

\begin{itemize}
\item \textbf{Douay-Rheims Bible (Challoner Revision).} Public-domain English text (Project Gutenberg eBook~1581). Cited at Exod.~3:14; Ps.~99:3, 103:29--30; Wis.~11:25--27; Matt.~4:10; Acts 17:27--28; Rom.~1:20, 11:33--36; 1~Cor.~4:7; Col.~1:17; Heb.~1:3.
\item \textbf{Vulgata Clementina.} Public-domain Latin text (eBible \texttt{latVUC} verse texts). Checked alongside the Douay-Rheims at Exod.~3:14; Ps.~99:3, 103:29--30; Wis.~11:25--27; Matt.~4:10; Acts 17:27--28; Rom.~1:20, 11:33--36; 1~Cor.~4:7; Col.~1:17; Heb.~1:3.
\end{itemize}

\subsection*{Ecumenical Councils and Magisterial Sources}

\begin{itemize}
\item \textbf{Fourth Lateran Council (1215).} Constitution~1, \emph{Firmiter credimus} (DS~800). Historical Latin text from Denzinger, \emph{Enchiridion Symbolorum} (public domain).
\item \textbf{First Vatican Council (1870).} Dogmatic Constitution on the Catholic Faith, \emph{Dei Filius}, ch.~1 (\emph{De Deo rerum omnium creatore}), ch.~2 (\emph{De revelatione}), ch.~4 (\emph{De fide et ratione}), can.~5 \emph{De Deo rerum omnium creatore}, and can.~1 \emph{De revelatione}. Historical Latin text and English translation by Cardinal Henry Edward Manning in \emph{Dogmatic Canons and Decrees} (New York: Devin-Adair, 1912), pp.~213--238 (public domain).
\item \textbf{Catholic Church.} \emph{Catechismus Catholicae Ecclesiae} (\emph{Catechism of the Catholic Church}), typical Latin edition (Vatican City: Libreria Editrice Vaticana, 1997), nn.~300--301. Cited by doctrinal paraphrase only (no verbatim reproduction).
\end{itemize}

\subsection*{Patristic Sources}

\begin{itemize}
\item \textbf{Athanasius of Alexandria, St.} \emph{Contra Gentes} (\emph{Against the Heathen}), \S\S41--42. Translated by Archibald Robertson. In \emph{A Select Library of Nicene and Post-Nicene Fathers of the Christian Church}, ser.~2, vol.~4 (New York: Christian Literature Company, 1892), pp.~26--27 (public domain).
\item \textbf{Augustine of Hippo, St.} \emph{Confessions}, Book~III, ch.~6 (11); Book~VII, ch.~10 (16), ch.~11 (17), ch.~17 (23). Translated by J.\ G.\ Pilkington. In \emph{Nicene and Post-Nicene Fathers}, ser.~1, vol.~1 (Buffalo: Christian Literature Company, 1886; public domain).
\item \textbf{Augustine of Hippo, St.} \emph{De Civitate Dei} (\emph{The City of God}), Book~X, ch.~4 (cited via Aquinas); Book~X, ch.~31 (cited via Aquinas); Book~XI, ch.~10, 24; Book~XII, ch.~1, 2, 6, 7, 25. Translated by Marcus Dods, vol.~1 (Edinburgh: T.\ \& T.\ Clark, 1871; Project Gutenberg eBook~45304, public domain).
\item \textbf{Augustine of Hippo, St.} \emph{De Genesi ad litteram}, Book~IV, ch.~12 (22--23). Latin text in J.-P.\ Migne, \emph{Patrologia Latina}, vol.~34 (Paris, 1845; public domain).
\item \textbf{Gregory of Nazianzus, St.} \emph{Oration}~28 (\emph{Second Theological Oration}), \S\S4--6, 17, 21; \emph{Oration}~38 (\emph{On the Theophany}), \S\S7--8. Translated by Charles Gordon Browne and James Edward Swallow. In \emph{Nicene and Post-Nicene Fathers}, ser.~2, vol.~7 (New York: Christian Literature Company, 1894; public domain).
\item \textbf{Irenaeus of Lyons, St.} \emph{Adversus Haereses} (\emph{Against Heresies}), Book~II, ch.~34 (3--4); Book~IV, ch.~39 (2). Translated by Alexander Roberts and W.\ H.\ Rambaut. In \emph{The Ante-Nicene Fathers}, vol.~1 (Buffalo: Christian Literature Publishing Co., 1887; public domain).
\item \textbf{John of Damascus, St.} \emph{De Fide Orthodoxa} (\emph{An Exact Exposition of the Orthodox Faith}), Book~I, ch.~9. Translated by S.\ D.\ F.\ Salmond in \emph{Nicene and Post-Nicene Fathers}, ser.~2, vol.~9 (Edinburgh: T.\ \& T.\ Clark, 1898; public domain); compared with the English Dominican rendering of Aquinas's Latin quotation in \emph{ST} I, q.~13, a.~11, corpus.
\item \textbf{Pseudo-Dionysius the Areopagite} (anonymous late fifth- or early sixth-century Christian author). \emph{De Caelesti Hierarchia} (\emph{The Celestial Hierarchy}), ch.~2 (\S3) and ch.~4 (\S1). Translated by John Parker (London: James Parker, 1899; public domain).
\end{itemize}

\subsection*{Medieval and Philosophical Sources}

\begin{itemize}
\item \textbf{Aristotle.} \emph{Physics}; \emph{Metaphysics}, Books~II, V, IX, XII. Analyzed in philosophical exposition (and quoted at \emph{Metaphysics} II, text~4 via Aquinas, \emph{ST} I, q.~44, a.~1).
\item \textbf{Duns Scotus, Bl.\ John.} \emph{Ordinatio}, Book~I, dist.~3 and dist.~8; Book~IV, dist.~11, q.~3. Referenced for the historic Catholic scholastic debate over univocity of being, intrinsic modes, and the distinction between essence and existence.
\item \textbf{Henry of Ghent.} \emph{Quodlibet}~I, q.~9 (1276). Referenced for the late thirteenth-century debate over an intentional distinction between essence and existence.
\item \textbf{Hume, David.} \emph{Dialogues Concerning Natural Religion}, Part~IX (London, 1779). Referenced in exposition of brute-fact objections to cosmological explanation.
\item \textbf{Plato.} \emph{Phaedo}; \emph{Republic}; \emph{Timaeus}. Analyzed in exposition of participation, the Form of the Good, and the limits of the Platonic Demiurge.
\item \textbf{Russell, Bertrand.} ``Debate on the Existence of God'' (BBC Third Programme broadcast with Frederick Copleston, S.J., 1948). Referenced in exposition of the brute-fact thesis.
\item \textbf{Su\'{a}rez, Francisco, S.J.} \emph{Disputationes Metaphysicae}, disp.~31 (Salamanca, 1597). Referenced for the early-modern Catholic scholastic alternative of a distinction of reason with a foundation in reality between essence and existence.
\item \textbf{Thomas Aquinas, St.} \emph{De Ente et Essentia}, ch.~4--5. Edited by L.\ Baur, emended by J.\ Koch (M\"unster: Aschendorff, 1933; Corpus Thomisticum Latin text, public domain).
\item \textbf{Thomas Aquinas, St.} \emph{Summa Theologiae}, Prima Pars (I, qq.~2, 3, 8, 12, 13, 19, 20, 21, 32, 44, 45, 46, 104, 105), Prima Secundae (I-II, q.~79, a.~2), and Secunda Secundae (II-II, q.~84, a.~1; q.~129, aa.~1, 3; q.~161, aa.~1--3). Translated by the Fathers of the English Dominican Province (New York: Benziger Brothers; Project Gutenberg eBooks 17611, 17897, and 18755, public domain).
\end{itemize}
